\documentclass[10pt,letterpaper]{article}

\usepackage[top=0.85in,left=2.75in,footskip=0.75in]{geometry}

% amsmath and amssymb packages, useful for mathematical formulas and symbols
\usepackage{amsmath,amssymb}

% Use adjustwidth environment to exceed column width (see example table in text)
\usepackage{changepage}

% Use Unicode characters when possible
\usepackage[utf8x]{inputenc}

% textcomp package and marvosym package for additional characters
\usepackage{textcomp,marvosym}

% cite package, to clean up citations in the main text. Do not remove.
\usepackage{cite}

% Use nameref to cite supporting information files (see Supporting Information section for more info)
\usepackage{nameref,hyperref}

% line numbers
\usepackage[right]{lineno}

% ligatures disabled
\usepackage{microtype}
\DisableLigatures[f]{encoding = *, family = * }

% color can be used to apply background shading to table cells only
\usepackage[table]{xcolor}

% array package and thick rules for tables
\usepackage{array}

% enumerate package lets us use letters instead of numbers
\usepackage{enumerate}

% create "+" rule type for thick vertical lines
\newcolumntype{+}{!{\vrule width 2pt}}

% create \thickcline for thick horizontal lines of variable length
\newlength\savedwidth
\newcommand\thickcline[1]{%
  \noalign{\global\savedwidth\arrayrulewidth\global\arrayrulewidth 2pt}%
  \cline{#1}%
  \noalign{\vskip\arrayrulewidth}%
  \noalign{\global\arrayrulewidth\savedwidth}%
}

% \thickhline command for thick horizontal lines that span the table
\newcommand\thickhline{\noalign{\global\savedwidth\arrayrulewidth\global\arrayrulewidth 2pt}%
\hline
\noalign{\global\arrayrulewidth\savedwidth}}

\usepackage{color}

% Remove comment for double spacing
%\usepackage{setspace}
%\doublespacing

% Text layout
\raggedright
\setlength{\parindent}{0.5cm}
\textwidth 5.25in
\textheight 8.75in

% Bold the 'Figure #' in the caption and separate it from the title/caption with a period
% Captions will be left justified
\usepackage[aboveskip=1pt,labelfont=bf,labelsep=period,justification=raggedright,singlelinecheck=off]{caption}
\renewcommand{\figurename}{Fig}

% Remove brackets from numbering in List of References
\makeatletter
\renewcommand{\@biblabel}[1]{\quad#1.}
\makeatother

% Leave date blank
\date{}

% Header and Footer with logo
\usepackage{lastpage,fancyhdr,graphicx}
\usepackage{epstopdf}
\pagestyle{myheadings}
\pagestyle{fancy}
\fancyhf{}
\setlength{\headheight}{27.023pt}
\lhead{\includegraphics[width=2.0in]{PLOS-submission.eps}}
\rfoot{\thepage/\pageref{LastPage}}
\renewcommand{\footrule}{\hrule height 2pt \vspace{2mm}}
\fancyheadoffset[L]{2.25in}
\fancyfootoffset[L]{2.25in}
\lfoot{\sf PLOS}


\begin{document}
\vspace*{0.2in}

\begin{flushleft}
{\Large
\textbf\newline{Ten Simple Rules for Writing a Technical Book}
}
\newline
\\
{Jess Haberman}\textsuperscript{1},
{Greg Wilson}\textsuperscript{2*}
\\
\bigskip
\textbf{1} Anaconda Inc., Austin, Texas, United States, 0000-0002-3779-1999\\
\textbf{2} Third Bit, Toronto, Ontario, Canada, 0000-0001-8659-8979\\
\textbf{*} gvwilson@third-bit.com.
\end{flushleft}

\section*{Introduction}

The Internet has put a world of information just a click away, but
that information is often fragmentary, contradictory, incomplete, or
out of date.  People still need coherent narratives that provide a
coherent overview and lay out a learning path that will get them
there.  In other words, they still need book-length expositions, which
means someone has to write them and publish them.

One of us (Haberman) has been helping writers and experts produce
instructional content for most of her career: prior to becoming
Director of Learning Solutions at Anaconda, she spent 14 years in book
publishing, with several years at O'Reilly Media signing and
developing technical books.  The other (Wilson) has written four books
on programming, co-authored three others, and edited six.  The tips
below summarize what we've learned along the way.

\section*{Rule 1: Don't write for money.}

The first thing any prospective author needs to understand is the
economics of publishing.  Typically, one-third of the books on a
publisher's list lose money, one-third break even, and one-third are
profitable to a greater or lesser extent.  For most, it's ``lesser'':
publishers actually make most of their revenue from a handful of
bestsellers.

For a technical book, your publisher will typically give you royalties
ranging from 10\% to 15\% of net revenue (i.e., revenue after
discounts).  You can expect to sell 200–2000 copies; a handful of
books do much better than that each year, but only a handful.
Downloading has obviously impact this: PDFs and BitTorrent files often
appear online within hours of the first digital sales of books
(particularly textbooks).  Assuming a US\$50 price tag, this means
that you will probably earn between US\$500 and US\$7500.  Since it's
going to take at least a thousand hours to write the book, you'll
almost certainly be better off financially doing some consulting (or
even getting a part-time job in a fast food restaurant).

You can earn much higher royalties by self-publishing, but doing so
comes with many tradeoffs.  You won't have access to the expertise of
a publisher when it comes to layout and design, cover art, several
levels of editing (development, copyediting, proofreading), printing,
print and digital distribution, marketing, and sales.  You must
dedicate time, money, and energy to finding partners for these tasks
or handle them yourself, so choose your path wisely.

So why write a book?

\begin{itemize}
\item
  \textbf{To build your reputation.}  Many musicians now view albums
  as a way to get people to buy concert tickets and t-shirts.
  Similarly, many authors and publishers view books as a way to get
  people to enrol in tie-in workshops and speaking engagements.
\item
  \textbf{To give back to the community.}  None of us got here without
  help, and we have a moral obligation to repay that by helping
  others.
\item
  \textbf{Because you enjoy doing it.}  Some people like working with
  wood or wool; you might like working with words.
\item
  \textbf{To gain a deeper understanding of the topic.}  Writing
  requires research, and you'll gain expertise during the writing
  process.
\end{itemize}

\section*{Rule 2: Read, then write.}

Go through two or three of your favorite technical books and make
notes about what their authors have done well.  Doing this will give
you some ideas for your own book, but it will also help you become
more conscious of your writing—in musical terms, it will help you
learn how to listen to yourself while you're playing.  One of us
(Wilson) learned a lot from the classic Unix books by Kernighan et al
\cite{Ke1979,Ke1981,Ke1983,Ke1988} and from \cite{Ud1999}, but your
tastes may vary.

Reading first will also help you understand what knowledge is already
public or well understood, which will help make your arguments even
stronger.

\section*{Rule 3: Test your material.}

Writing a book is a big undertaking, and it's easy to get discouraged.
Instead of tackling it head-on, start blogging regularly (e.g., one
article a week, each a few hundred to a thousand words long).  Don't
try to write them in the order you think they'll eventually go into
the book: one reason to blog is to get a better idea of what should
come before what.  As you become more comfortable, though, try writing
multi-post series that can eventually become chapters.

If you run even a single workshop based on the material for your book
while you are writing it, you will almost certainly realize that
several of your core assumptions about your audience and material are
completely wrong \emph{while you still have time and energy to fix
things}.  Conferences are always looking for tutorials, your local
college or university might be looking for a sessional course, and
dozens of learning sites exist to publish training materials (though
in this case you must be careful that their licensing allows you to
re-use the material).  Your colleagues will probably welcome some
lunch \& learn sessions, and if all else fails, you can create a
YouTube channel to find out how your material sounds.

A corollary is that there's no point trying material out if you don't
\textbf{incorporate what you learn}.  If you pay attention to
feedback, you will probably rewrite every line of your book at least
twice and revise some parts a dozen times or more.  You can skip this
work if you don't collect or act on feedback, but your book will
almost certainly be a lot worse for it.

We recommend that you drywall and then paint, i.e., that you write a
rough draft with placeholders and \emph{then} worry about diagrams,
bibliography citations, glossary entries, and so on.  For example, the
word \texttt{FIXME} appeared over a hundred times in first
``complete'' draft of \cite{Wi2019}; half of those turned out to be
irrelevant by the time the text settled down.

An important special case of this is that you should do diagrams on a
whiteboard first.  It's faster and more flexible than any computer
drawing tool (even fingertip sketching apps for tablets) and a photo
taken with a cell phone is good enough for your first readers.  Many
publishers will re-create your sketches anyway, so don't fuss over
them too much.

\begin{mdframed}
\textbf{Union and Intersection}

\noindent
No matter how carefully you plan, you will hit part 5 and realize that
there was something you should have explained back in part 3.  You can
fill this in on the fly if you are teaching in person, but that
doesn't work in a book.  For this reason, a live lesson should be the
\emph{intersection} of what audience members don't know (i.e., the
things you're sure they all need) while a book has to be the
\emph{union} of what readers don't know (i.e., the things any of them
need).
\end{mdframed}

\section*{Rule 4: Start with a learner persona.}

A \emph{learner persona} is a fictional character that captures key
aspects of who you are trying to teach, what they already know, and
what they want to learn \cite{Wi2019}.  Create one or two that are
fairly similar (i.e., a primary and secondary audience) and write for
them, because a book that's meant for everyone is actually useful to
no-one.  Imagine what your persona already knows, their career or
academic level, and what challenges they face.  Your goal is to build
on the former to offer solutions to the latter.

As you're doing this, keep in mind the difference between a
\emph{novice} (who is trying to build a mental model), a
\emph{competent practitioner} (who has one and wants to fill in gaps
in their knowledge), and an \emph{expert} (who wants higher-level
discussion of tradeoffs and alternatives).  Novices need tutorials
that help them build a coherent mental model of the domain; competent
practitioners want to fill in the gaps in their knowledge, and experts
are usually interested in higher-level discussions of patterns,
alternatives, and tradeoffs.  No single book can serve all three well.

As a corollary, \textbf{don't write a general introduction to a topic}
unless you're going to be one of the first two or three to market.
Instead, focus on a particular aspect (e.g., ``Python for Web
Scraping'') or a particular audience (e.g., ``R for Librarians'').
Your potential audience may be smaller, but you'll reach a much higher
percentage of that audience, and a focused book is much easier to
write.

And please, \textbf{don't write a reference book}: it can never be as
comprehensive or stay as up-to-date as the Internet.

\section*{Rule 5: Know your competition.}

Conduct research to find out what other resources exist for this
topic, keeping in mind they don't have to be books.  Where is the
audience that you have already identified learning more about the
topic right now?  What can you offer that differs from and improves on
that?  The most common weakness in other resources is they are out of
date or no longer accurate, but other times they are poorly organized,
not well edited, lack some unique insight you can bring to the topic,
or are trying to be too general.  Every publisher's proposal form will
ask about competition; finding an answer \emph{before} you have
written several chapters can save you a lot of fruitless labor.

Note that some differentiators are not significant enough to be
selling points.  For example, you might argue that your book is better
because it has color diagrams while competing titles are black and
white.  Color diagrams usually don't make a significant difference in
comprehension, so to an editor, a request for color printing usually
equates to more investment and less profit.

\section*{Rule 6: Avoid common mistakes.}

\begin{itemize}
\item
  \textbf{Avoid banal advice} (like ``avoid common mistakes'').  Few
  things are as frustrating to the reader as sentences like, ``You
  should carefully consider users' needs,'' because no sensible person
  would recommend the opposite.  If you can't provide a checklist of
  specific things to consider, or a scenario to give the reader
  insight into how you think through that class of problem, you're
  giving them the intellectual equivalent of junk calories.
\item
  \textbf{Don't be too granular.}  Don't spend months or years writing
  a book in such specific detail that a new release of software makes
  your advice obsolete or simply wrong.  Instead, generalize and focus
  on principles that are likely to be long-lived without violating the
  rule about banality.  Your goal should be for your book to remain
  relevant for three to five years.
\item
  \textbf{Don't write to cover your ass.}  We have read dozens of
  books that start with a short introduction to XYZ.  We cannot
  remember a single one that was actually useful.  If your audience
  knows Python, they don't need a chapter-length intro to the
  language, and if they don't, one isn't going to help them.  Authors
  usually include those short introductions so that they can claim the
  book is suitable for novices when they knew in their hearts it
  isn't.  Remember, it's perfectly acceptable to target an audience
  with more sophisticated knowledge than ``none'' or ``basic''.
\item
  \textbf{Don't try to write.}  Don't try to sound like an expert;
  instead, lean into the idea that you \emph{are} an expert and write
  in your own voice.  You don't convince readers (or publishers) you
  have worthwhile opinions by sounding like what you imagine smart
  people sound like.
\item
  \textbf{Don't try to be funny.}  Very few jokes are funny the second
  time you hear them; even fewer can stand a third re-telling unless
  they're very, very dry.  On a related note, please don't use
  exclamation marks: what you're writing probably isn't that
  surprising, and certainly won't be the second time around.  Most
  bold, italics, and underline formatting is superfluous\footnote{And
  remember: nobody reads footnotes.}.
\end{itemize}

\section*{Rule 7: Write a proposal.}

If your goal is to have your book professionally published, write a
proposal before writing the book.  You should spend a few weeks on
this and make it as complete and convincing as possible; if nothing
else, it will help you figure out exactly what you're writing about.

When you pitch it to your preferred publishers, do some research to
find the likely editors so you can be sure they see it instead of it
winding up in the slush pile of unsolicited proposals and manuscripts.

\begin{mdframed}
\textbf{Shop Around}

\noindent
Know what your publishing goals are.  Do you prefer the style, brand,
or reach of a specific publisher?  Are you seeking the top financial
offer?  Identify your goals before you begin pitching.  Don't feel you
have to send your proposal to one publisher at a time, but if you are
shopping it around, it is common courtesy to let editors know if the
proposal has garnered interest elsewhere.  Communicating your goals
and timelines can only help you in navigating the pitch process
smoothly and achieve your publishing goals.
\end{mdframed}

If you have a strong proposal, it should at the very least land you
some time with editors who are knowledgable about the market for the
topic.  The best editors will help you hone the proposal by asking
good questions and steering your book in the right direction.

If you receive an offer for publication, negotiate.  Let editors know
if you have received or anticipate an offer elsewhere, or if you're
considering self-publishing.  There's no harm in seeking the best
offer you can (i.e., the one most aligned with your goals, whatever
they are): even if an offer doesn't improve, you've wasted nothing and
can move forward knowing that others are truly invested in your book
and will help keep you accountable.

\section*{Rule 8: Automate, but proofread.}

This is the technical author's equivalent of ``trust, but verify''.
Depending on the tools you use, you may be able to regenerate all of
the examples in your book with a single command, but you must still
re-read the discussion about those examples every time you make a
change to make sure they haven't fallen out of step.  Similarly, it
doesn't matter how well you know your chosen subject: something will
have changed since the last time you looked, and something else will
change in the year or more it takes you to complete your manuscript.
Your publisher will keep track of errata (i.e., printing errors); it
improves longevity of your book if you address them at regular
intervals.

Keep in mind that most people won't read your book linearly: they will
skip some sections or jump ahead to find an answer to a problem they
need to solve right away.  You can help them by making liberal use of
cross-references (e.g., ``as covered in Chapter 2'') to help them
navigate \cite{Li2020}.  And because page and chapter numbers
frequently change, use chapter titles instead.

Speaking of chapter titles, consider titles that begin with verbs.  If
you're writing an instructional book, this will help ensure you're
focused on the task or skill you're trying to teach.  (See
\cite{Me2022} for an example.)

\section*{Rule 9: Trust your reviewers and your editor. }

Lots of people will do cursory reviews or give your manuscript a
single careful read; people who will look at changes over and over
again and see what's actually on the page each time are worth their
weight in rubies.  \emph{Trust them}: if a passage doesn't make sense
to them, then it doesn't matter if it makes sense to you.  Find a
colleague or friend who you can trust to be brutally honest.  It's one
thing to get feedback from your editor or outside reviewers, but
someone you trust who will challenge your arguments is invaluable.

Similarly, as Tom Wilkie once said, an author's job is to produce the
manure in which an editor grows something worth reading.  You may
sweat over an example for days, but your audience doesn't see your
effort: they see your results.  If explanations that seem clear,
concise, and elegant to you don't make sense to your readers, then
they're right and you're wrong.

A professional editor doesn't get paid extra for making more edits. So
if they mark something, they probably have a good reason. Ask why if
you're not sure. Keep in mind that your motivation and your editor's
are the same, to make the book the best it can possibly be.

Don't be afraid to set it aside. One of us (Wilson) wrote three
quarters of an introduction to R for Python programmers that will
probably never see the light of day.  This is not unusual: many
authors find that the more they write, the less they believe the book
will actually help its intended audience.  There is nothing to be
ashamed of in this: scientists don't expect every hypothesis to turn
out to be true, and authors shouldn't either.

Of course, if you have signed a contract and are suddenly rethinking
whether your book should exist, have an honest conversation with your
editor about next steps.  Don't hide information, and don't become
impossible to reach.  Editors are used to authors becoming discouraged
and can provide a much-needed reality check.  They can also help you
evolve the reasoning, scope, timing, or approach to your book: it's
actually part of their job to do so.

\section*{Rule 10: Stop, ship, and celebrate.}

Your book will never be perfect.  It will probably never even feel
\emph{finished}, any more than software does, but you should ship it
anyway and then celebrate what you've accomplished.  Writing a book
can be a slog, but it's not easy to replicate the feeling of
accomplishment when you find your book listed online or receive a
shipment of printed copies.  Don't be afraid to tell your friends or
post on your social media accounts: it isn't bragging if you actually
did it, and if it was worth your time to write the book, it's worth
your time to help people find it.

\bibliographystyle{unsrt}
\bibliography{11-techbook}

\end{document}
